\begin{tikzpicture}[every node/.append style={execute at begin node=\setstretch{1.05}}, font=\small, x=1cm, y=1cm]
  % 1024 enderecos em 15 cm
  \fill[vvisi!18] (0,0) rectangle (7.5,1.0);
  \fill[valun!18] (7.5,0) rectangle (11.25,1.0);
  \fill[vprof!28] (11.25,0) rectangle (12.1875,1.0);
  \fill[black!5] (12.1875,0) rectangle (15,1.0);
  \draw (0,0) rectangle (15,1.0);
  \foreach \x in {7.5, 11.25, 12.1875, 13.125} \draw (\x,0) -- (\x,1.0);
  \draw[densely dotted, black!45] (3.75,0) -- (3.75,1.0);
  \node[text=black!65] at (1.875,0.5) {192.168.0.x};
  \node[text=black!65] at (5.625,0.5) {192.168.1.x};
  \node[text=black!65] at (9.375,0.5) {192.168.2.x};
  \node[text=black!55, font=\footnotesize] at (12.656,0.5) {/26};
  \node[text=black!55, font=\footnotesize] at (14.06,0.5) {/25};
  \draw[decorate, decoration={brace, amplitude=5pt}] (0,1.1) -- node[above=5pt, align=center] {\textbf{Visita, VLAN 30}\\192.168.0.0/23: 512 endereços} (7.5,1.1);
  \draw[decorate, decoration={brace, amplitude=5pt}] (7.5,1.1) -- node[above=5pt, align=center] {\textbf{Alunos, VLAN 20}\\192.168.2.0/24: 256} (11.25,1.1);
  \draw[decorate, decoration={brace, amplitude=4pt}] (12.1875,1.1) -- node[above=5pt, align=center, text=black!60] {livre\\192} (15,1.1);
  \draw (11.72,1.0) -- (11.72,2.55);
  \node[anchor=south, align=center] at (11.72,2.55) {\textbf{Professor, VLAN 10}\\192.168.3.0/26: 64};
  \foreach \x/\t in {0/.0.0, 3.75/.1.0, 7.5/.2.0, 11.25/.3.0, 12.1875/.3.64, 13.125/.3.128, 15/.3.255} {
    \draw (\x,0) -- (\x,-0.15);
    \node[below, font=\footnotesize] at (\x,-0.15) {\t};
  }
  \node[anchor=north west, font=\footnotesize, text=black!60, text width=15cm] at (0,-0.65) {As marcas mostram o fim do endereço (192.168 omitido). A linha pontilhada separa 192.168.0.x de 192.168.1.x, que pertencem à mesma subrede /23 da Visita.};
\end{tikzpicture}
